\singlespacing
In Local Search, the core hyperparameters include the number of iterations, the neighborhood sample size, and the swap probability. The sample neighbors parameter dictates the breadth of the algorithm's vision at each step; a larger sample size increases the likelihood of finding a steeper descent path, potentially leading to a lower peak memory. Concurrently, however, the execution time grows in linear proportion to the product of iters and sample neighbors. Furthermore, due to the algorithm's greedy nature of rejecting inferior solutions, it tends to stagnate early once a local optimum is reached. In such cases, setting an excessively large iters does not yield further reductions in peak memory and only results in futile computational overhead.

For Simulated Annealing, the final memory optimization outcome is highly dependent on the cooling schedule, governed by the initial temperature, the cooling rate, and the total iterations. A slower cooling rate allows the algorithm to maintain a higher probability of accepting inferior solutions for a longer period, effectively enabling it to escape local optima traps and discover a lower peak memory. However, to accommodate this gradual cooling process and ensure the convergence of the Markov chain, the iters parameter must be drastically increased. Because the execution time of Simulated Annealing is directly proportional to iters, pursuing extreme memory compression inherently necessitates a significant increase in computation time.

In the Genetic Algorithm, the population size, number of generations, and mutation rate collectively determine the depth and breadth of the search. Expanding the pop size introduces greater diversity into the candidate solutions, while increasing generations allows superior genetic traits to be more thoroughly crossed and evolved. Elevating both parameters significantly improves the stability and success rate of converging to the global optimum. The mutation rate serves as a perturbation mechanism to prevent premature convergence of the population. Regarding computation time, the total execution delay is proportional to the product of pop size and generations. Consequently, by judiciously tuning these two parameters, one can precisely control the trade-off between extensively exploring the solution space and minimizing the compilation delay.
